\section{Conclusions}\label{s:conclusion}
% \paragraph{Summary.}
This work presents a spectral shape–based framework for large-scale analysis of LLMs that is data-free, computationally efficient, and robust to post-training transformations. We curate and benchmark a large and diverse LLM corpus and conduct extensive experiments across similarity measurement, classifications, clustering, and performance prediction. Our results have shown the consistent effectiveness of our spectral signature across various experimental setups. Furthermore, we validate that ESD shape metrics, particularly \ALPHAHILL, capture intrinsic structural properties of LLMs that are stable under output-invariant reparameterizations and noise perturbations. Therefore, our spectral signature can provide a practical tool for managing modern LLM ecosystems, with implications for model provenance, organization, and scalable evaluation.

% \paragraph{Limitations and Future Work.}
% Spectral shape signatures provide a coarse and global characterization and may not capture task-specific behaviors or fine-grained functional differences. Future work will explore extending the framework to broader model classes (e.g., encoder–decoder and multimodal models), integrating layer- or module-specific signatures for finer resolution, and combining spectral descriptors with lightweight data-aware signals to improve sensitivity while preserving scalability and efficiency.